\documentclass[11pt]{article}

\usepackage[margin=1in]{geometry}
\usepackage{amsmath, amssymb, amsthm}
\usepackage{mathtools}
\usepackage[colorlinks=true, linkcolor=blue, urlcolor=blue, citecolor=blue]{hyperref}

\title{\textbf{The spectral theorem, intuitively}}
\author{Sanjay Suresh}
\date{February 2026}

\begin{document}
\maketitle

The spectral theorem says every real symmetric (or complex Hermitian) matrix can be written as
\[ A = Q\,\Lambda\,Q^{\mathsf{T}}, \]
an orthonormal basis $Q$ of eigenvectors and a diagonal matrix $\Lambda$ of real eigenvalues. The
geometric reading is the useful one: such a matrix acts as pure stretching along a set of
perpendicular axes, with no rotation or shear.

That is why these matrices are everywhere I work. A Hamiltonian is Hermitian, so its eigenvalues are
real energies and its eigenvectors are stationary states. A covariance matrix is symmetric, so PCA
is just reading off its eigen-axes. Whenever a problem reduces to a symmetric operator, the spectral
theorem promises a clean coordinate system in which it becomes diagonal.

\end{document}
